\documentclass[oneside,a4paper,14pt]{extarticle} %размер шрифта 14
\usepackage[T1,T2A]{fontenc}
\usepackage[
a4paper,
letterpaper,
top=2cm,
bottom=2cm,
left=2.5cm,
right=1.5cm
]{geometry}
\usepackage[utf8]{inputenc} %кодировка текста
\usepackage[russian]{babel} %поддержка русского языка
\usepackage{textcomp} %текстовые символы
\usepackage{indentfirst} %корректировка отступов
\usepackage{listings} %для кода !!!
\usepackage{listingsutf8}
\usepackage{xcolor}
\usepackage{float}
\usepackage{enumitem} %для нумерации буквами
\usepackage{graphicx} %работа с изображениям
\usepackage{mwe} % for blindtext and example-image-a in example
\usepackage{wrapfig}
\usepackage{caption}
\usepackage{amsmath}  % для формул и символов
\usepackage{amsfonts}
\usepackage{amsthm}
\usepackage{graphicx}
\usepackage[all]{xy}
\usepackage[breaklinks]{hyperref}
\usepackage{xurl}
% размер кегля у заголоков разделов
\usepackage{titlesec}

\titleformat{\section}
{\normalfont\large\bfseries}{\thesection}{1em}{}
\titleformat{\subsection}
{\normalfont\normalsize\bfseries}{\thesubsection}{1em}{}
\titleformat{\subsubsection}
{\normalfont\normalsize\bfseries}{\thesubsection}{1em}{}
\renewcommand\baselinestretch{1.33}
\normalsize % межстрочный интервал
\setlength{\parindent}{1.25cm}
\usepackage{indentfirst}

% настройки листинга
\lstset{
  language=Python,
  basicstyle=\ttfamily\small,
  keywordstyle=\color{blue},
  commentstyle=\color{green!60!black},
  stringstyle=\color{red},
  numbers=left,
  numberstyle=\tiny\color{gray},
  frame=single,
  breaklines=true,
  showstringspaces=false,
  tabsize=2,
  inputencoding=cp1251,
  extendedchars=true,
}

\begin{document}
	\newpage\thispagestyle{empty}
	\begin{center}
		МИНИСТЕРСТВО НАУКИ И ВЫСШЕГО ОБРАЗОВАНИЯ\\
			РОССИЙСКОЙ ФЕДЕРАЦИИ
			ФЕДЕРАЛЬНОЕ ГОСУДАРСТВЕННОЕ БЮДЖЕТНОЕ\\
			ОБРАЗОВАТЕЛЬНОЕ
			УЧРЕЖДЕНИЕ ВЫСШЕГО ОБРАЗОВАНИЯ\\
			«ВЯТСКИЙ ГОСУДАРСТВЕННЫЙ УНИВЕРСИТЕТ»\\
			Институт математики и информационных систем\\
			Факультет автоматики и вычислительной техники\\
			Кафедра электронных вычислительных машин
	\end{center}
	\vspace{20mm}

	\begin{center}
		Отчёт по лабораторной работе №7\\
		по дисциплине\\
		<<Промпт-инжиниринг>>\\
		<<Разработка инструментов автоматизации на базе API больших языковых моделей>>\\
		<<Вариант 4>>\\
	\end{center}
	\vspace{48mm}

	\noindent
	Выполнил студент гр. ИВТб-1303-06-00 \hfill \rule[-0,5mm]{30mm}{0.15mm}\,/Гортоломей И.К./
	\newline
	Проверил доцент кафедры ЭВМ \hfill \rule[-0,5mm]{30mm}{0.15mm}\,/Колупаев А.В./

	\vfill
	\begin{center}
		Киров\\
		2026
	\end{center}

	\newpage
	\thispagestyle{empty}

\section*{Введение}
\noindent
\textbf{Цель работы:} Формирование профессиональных компетенций в области проектирования и реализации программных систем, использующих возможности больших языковых моделей (LLM) через API. Необходимо научиться инкапсулировать логику промпт-инжиниринга в программный код, создавать устойчивые к ошибкам алгоритмы генерации контента и освоить методы автоматической валидации и структурирования выходных данных нейросетей для решения прикладных инженерных задач.

\noindent
\textbf{Используемые модели и API:}
\begin{itemize}
    \item OpenRouter API с моделью \texttt{google/gemma-4-31b-it:free} --- основная LLM для генерации учебных программ.
    \item Python 3.10+ --- язык реализации.
    \item Библиотеки: \texttt{openai}, \texttt{json}, \texttt{time}, \texttt{google.colab.userdata} (для безопасного хранения ключа).
\end{itemize}

\noindent
\textbf{Предметная область варианта 4: Конструктор учебных программ} \\
Программа автоматически создаёт детальные учебные программы по заданной дисциплине и уровню подготовки (Школа/ВУЗ/Курсы). Результат включает: почасовой план занятий, список необходимой литературы, темы для домашних заданий и контрольных работ. Всё адаптируется под указанный уровень подготовки. Результат сохраняется в JSON и форматированном текстовом документе.

\newpage
\section{Задание №1: Проектирование архитектуры генератора промптов}
\subsection{Анализ требований варианта}
\noindent
Вариант №4 (тип «Генератор текстов») требует:
\begin{itemize}
    \item \textbf{Переменные:} Название дисциплины, \\ Уровень подготовки (Школа/ВУЗ/Курсы).
    \item \textbf{Константы:} Системная роль («Ты — опытный методист...»), структура учебной программы, формат вывода.
    \item \textbf{Этапы:} 1) Сгенерировать общую структуру курса; 2) Детализировать почасовой план; 3) Подобрать литературу; 4) Создать темы ДЗ.
\end{itemize}

\subsection{Интерфейс класса PromptBuilder}
\noindent
Создан класс, отвечающий только за сборку текста промпта (не обращается к API).
\begin{itemize}
    \item \texttt{\_\_init\_\_()} --- инициализация системной роли.
    \item \texttt{set\_discipline(discipline)} --- установка названия дисциплины.
    \item \texttt{set\_level(level)} --- установка уровня подготовки.
    \item \texttt{build\_program\_prompt()} --- возвращает список сообщений для генерации учебной программы.
    \item \texttt{get\_system\_role()} --- геттер системной роли.
\end{itemize}

\newpage
\subsection{Листинг класса PromptBuilder}
\lstinputlisting{1.nonutf.txt}

\newpage
\subsection{Тестирование сборки промптов (без API)}
\noindent
Для проверки класса был выполнен следующий тестовый код:
\lstinputlisting{2.nonutf.txt}

\noindent
\begin{figure}[H]
  \centering
  \includegraphics[width=17cm]{1.png}
  \caption{результат тестирования класса PromptBuilder}
\end{figure}

\newpage
\section{Задание №2: Реализация программного взаимодействия с API}
\subsection{Выбор API и параметры генерации}
\noindent
Использован OpenRouter API (бесплатный доступ) с моделью \\ \texttt{google/gemma-4-31b-it:free}.

\begin{table}[H]
\centering
\begin{tabular}{|l|l|p{11cm}|}
\hline
\textbf{Параметр} & \textbf{Значение} & \textbf{Обоснование} \\ \hline
temperature & 0.3 & Низкая температура для точного следования формату JSON и уменьшения галлюцинаций. \\ \hline
top\_p & 0.95 & Стандартное значение, обеспечивающее разумное разнообразие. \\ \hline
\end{tabular}
\end{table}

\noindent
API-ключ хранится в секретах Google Colab \\ (переменная OPENROUTER\_API\_KEY), доступ через \texttt{userdata.get()}.

\subsection{Реализация основного конвейера}
\noindent
Скрипт выполняет следующие шаги:
\begin{enumerate}
    \item Создание экземпляра PromptBuilder с параметрами.
    \item Генерация учебной программы через \texttt{generate\_curriculum()}.
    \item Автоматическая очистка JSON от постороннего текста.
    \item Сохранение программы в \texttt{curriculum\_\{discipline\}.json} и \texttt{.txt}.
\end{enumerate}

\newpage
\subsection{Листинг функций взаимодействия с API}
\lstinputlisting{3.nonutf.txt}

\subsection{Пример выполнения и структуры JSON-ответа}
\begin{figure}[H]
  \centering
  \includegraphics[width=17cm]{2.png}
  \caption{результат выполнения программы}
\end{figure}
\begin{figure}[H]
  \centering
  \includegraphics[width=17cm]{4.png}
  \caption{Фрагмент текстового файла с отформатированной учебной программой}
\end{figure}
\newpage
\begin{figure}[H]
  \centering
  \includegraphics[width=15cm]{3.png}
  \caption{Фрагмент JSON-файла}
\end{figure}
\newpage

\section{Задание №3: Тестирование, оценка качества и универсальности}
\subsection{Тестовые сценарии}
\noindent
Программа запущена на трёх различных наборах входных данных:
\begin{enumerate}
    \item Дисциплина «Программирование на Python», уровень «Курсы»
    \item Дисциплина «История России», уровень «Школа»
    \item Дисциплина «Машинное обучение», уровень «ВУЗ»
\end{enumerate}

\noindent
\textbf{Пример «сырого» ответа модели до очистки:}
\begin{verbatim}
Хорошо, вот учебная программа по вашему запросу:

```json
{
  "discipline": "Программирование на Python",
  "level": "Курсы",
  ...
}
```
Надеюсь, эта программа будет полезна!
\end{verbatim}
\newpage

\subsection{Результаты тестирования}
\begin{table}[H]
\centering
\begin{tabular}{|l|p{3cm}|l|p{2cm}|l|l|}
\hline
\textbf{Тест} & \textbf{Дисциплина} & \textbf{Уровень} & \textbf{JSON валиден} & \textbf{Часы} & \textbf{Литература/ДЗ} \\ \hline
1 & Програм- мирование на Python & Курсы & Да & 36 часов & 2+1 / 3 темы \\ \hline
2 & История России & Школа & Да & 34 часа & 3+2 / 4 темы \\ \hline
3 & Машинное обучение & ВУЗ & Да & 72 часа & 4+3 / 5 тем \\ \hline
\end{tabular}
\end{table}

% \noindent
% МЕДИА: [Скриншоты трёх сгенерированных текстовых файлов с учебными программами для разных уровней]

\subsection{Анализ ошибок и их исправление}
\begin{table}[H]
\centering
\begin{tabular}{|l|p{5cm}|p{5cm}|p{5cm}|}
\hline
\textbf{№} & \textbf{Проблема} & \textbf{Корректировка} & \textbf{Результат} \\ \hline
1 & Модель возвращала JSON с лишним текстом & Добавлена очистка: поиск символов \texttt{\{} и \texttt{\}} с обрезкой строки & Парсинг стал стабильным \\ \hline
2 & Для школьного уровня программа была слишком сложной & Усилен системный промпт: явное указание «используй доступный язык» & Программы адаптированы под уровень \\ \hline
3 & Иногда не хватало токенов & Убран max\_tokens & Все разделы генерируются полностью \\ \hline
\end{tabular}
\end{table}

\noindent
\textbf{Вывод по тестированию:} Инструмент успешно генерирует учебные программы для различных дисциплин и уровней подготовки. Программа универсальна: достаточно изменить входные параметры, чтобы получить новую программу без изменения кода.

\newpage
\section*{Заключение}
\noindent
\textbf{Оценка эффективности автоматизации:} Разработанный конструктор учебных программ полностью автоматизирует создание детальных образовательных программ. Время обработки одного запроса составляет около 30-60 секунд. При ручном составлении аналогичного объёма потребовалось бы не менее 2-4 часов. Экономия времени: в 120-240 раз.

\noindent
\textbf{Выявленные ограничения:}
\begin{itemize}
    \item Лимиты токенов: для очень подробных программ может потребоваться увеличение max\_tokens.
    \item Скорость API: бесплатные тарифы OpenRouter имеют задержки.
    \item Качество русского языка: Gemma-4 иногда использует неестественные обороты.
    \item Актуальность литературы: требуется ручная проверка изданий.
\end{itemize}

\noindent
\textbf{Наиболее эффективные техники промпт-инжиниринга:}
\begin{itemize}
    \item Few-Shot через JSON-схему: добавление детальной структуры ожидаемого JSON улучшило качество ответов.
    \item Conditional Prompting: различные инструкции для разных уровней подготовки позволили адаптировать программы без изменения кода.
    \item Self-Correction: автоматическая очистка ответа от лишнего текста позволила избежать падений при парсинге.
    \item Разделение ответственности: класс PromptBuilder изолирует конструирование промптов.
\end{itemize}

\section*{Приложения}
\noindent
\textbf{Приложение А. Тестовые данные}
\lstinputlisting[language=]{5.nonutf.txt}

\noindent
\textbf{Приложение Б. Ссылки} \\
Google Colab Notebook: \url{https://colab.research.google.com/drive/1g8o3OJALW7QbAk1AI3AONqY8urqNdfj7?usp=sharing} \\
GitHub: \url{https://github.com/GIKExe/Subject/tree/main/Prompt/laba7}

\end{document}